\documentclass[10pt,a4paper]{amsart}
\usepackage[margin=19mm]{geometry}
\usepackage{amsmath,amssymb,mathtools}
\usepackage[scr=rsfs,cal=euler]{mathalfa}
\usepackage{xfrac}
\usepackage[unicode,hidelinks,bookmarks=false]{hyperref}
\newcommand{\ie}{i.e.\ }
\newcommand{\cf}{cf.\ }
\newcommand{\resp}{respectively\ }
\newcommand{\Subsection}{\subsection}
\providecommand{\nobreakdash}{\nobreak\hbox{-}}
\setlength{\emergencystretch}{2em}
\multlinegap0pt
\usepackage{bm}
\usepackage{bbm}
\usepackage{dsfont}
\usepackage{xspace}
\usepackage{cancel}
\usepackage{mathtools}
\usepackage{amsthm}
\makeatletter
\@ifundefined{thm}{\newtheorem{thm}{Theorem}}{}
\@ifundefined{lem}{\newtheorem{lem}[thm]{Lemma}}{}
\makeatother
\newcommand{\C}{\ensuremath{\mathbb{C}}}  
\newcommand{\N}{\ensuremath{\mathbb{N}}}
\newcommand{\Z}{\ensuremath{\mathbb{Z}}}
\newcommand{\cT}{\ensuremath{\mathcal{T}}}
\newcommand{\cU}{\ensuremath{\mathcal{U}}}
\title[A new source of purely finite matricial fields]{A new source of purely finite matricial fields}
\author{David Gao, Srivatsav Kunnawalkam Elayavalli, Aareyan Manzoor, Gregory Patchell}
\date{Source: arXiv:2603.24502v5 (CC BY 4.0)}
\begin{document}
\maketitle

\noindent\emph{Source and licence.} A new source of purely finite matricial fields. Version arXiv:2603.24502v5 (CC BY 4.0), distributed under the Creative Commons Attribution 4.0 International licence, as recorded in the arXiv licence metadata for this exact version and shown on the abstract page https://arxiv.org/abs/2603.24502v5. Full rights notice: licenses/arxiv-2603.24502-CC-BY-4.0.txt.\par\smallskip

\noindent\textit{Editorial scope.} Counted unit: the Lem whose source label is \texttt{lem:2.3} in the article (source0.tex lines 250--257, source label \texttt{lem:2.3}), delivered together with the article's own complete original proof of it (source0.tex lines 259--336, one \texttt{proof} environment of 78 lines). No mathematics is written, repaired or re-derived: statement and proof are the article's own text line for line. Required context retained uncounted: isomorphism thm (lines 232--234). Every definition, display and label of those lines is retained unchanged. Omitted from this package and not redistributed: 1--231, 235--249, 258--258, 337--437. Nothing inside a retained excerpt is omitted or altered: every retained body is delivered exactly as the article's lines are written, so no omission receipt of any kind accompanies this package. Citations: the retained excerpts carry no citation command at all, so no bibliography entry of the article is transcribed and no reference is redistributed. Reference adaptation: none was needed and none was made; every reference inside the delivered unit resolves inside it, so no unresolved reference or citation is produced. Printed slips kept literal: no printed word, symbol, hypothesis or number of the counted statement or of its proof was altered. Layout and class: the article's own class is \texttt{amsart} with packages including amsmath, amsthm, amsfonts, amssymb, enumerate, hyperref, xcolor, verbatim, comment; the wrapper is the lane's pinned \texttt{amsart} editorial wrapper at 10pt on A4, which restates the theorem-like environments the retained text needs and adds the article's own macro definitions that the retained text needs (\texttt{\textbackslash C}, \texttt{\textbackslash N}, \texttt{\textbackslash Z}, \texttt{\textbackslash cT}, \texttt{\textbackslash cU}), with extra wrapper packages (bm, bbm, dsfont, xspace, cancel, mathtools). The delivered numbering is the unit's own. Assets: no retained span contains a figure environment, a graphics-inclusion command, a PDF-page inclusion, a table or a TikZ picture; no image file is copied into this package, the package declares no asset dependency and no image file is redistributed with it. Licence: version arXiv:2603.24502v5 of this article is distributed under the Creative Commons Attribution 4.0 International licence, as recorded in the arXiv licence metadata for this exact version and shown on the abstract page https://arxiv.org/abs/2603.24502v5. A full rights notice is delivered at licenses/arxiv-2603.24502-CC-BY-4.0.txt. Characters: every retained excerpt is pure ASCII, so the delivered engine, XeLaTeX with its default Latin Modern text font, reports no missing character.
\begin{thm}\label{isomorphism thm}
    Let the notation be as above. Then the $C^*$-subalgebra generated by $A$ and $\frac{1}{2}(T+T^*)$ in $\cT(E)$ is isomorphic to $A\ast_{B} (B\otimes C[-1,1])$, where $\frac{1}{2}(T+T^*)$ is the semicircular in $C[-1,1]$.
\end{thm}

\begin{lem}\label{lem:2.3}
    Let $G$ be a group and $H$ a subgroup with $H=\bigcap_{i\in \N} H_i$ for some decreasing sequence of subgroups $H_i$. Let $K_i = \bigcap_{g \in G} gH_ig^{-1}$ be the normal core of $H_i$. Let $\cU$ be a free ultrafilter on $\N$. Then there is a trace-preserving embedding of $C^*$-algebras:
    \[C_r^*(G\ast_H (H\times\Z)) \hookrightarrow \prod_\cU C_r^*([G/K_i\ast_{H_i/K_i} (H_i/K_i\times\Z)] \times G).\]

    Moreover, this embedding lifts to a sequence of group homomorphisms
    \[G \ast_H (H \times \Z) \to [G/K_i\ast_{H_i/K_i} (H_i/K_i\times\Z)] \times G\] which sends $g\in G$ to $(gK_i, g)$ and $n\in \Z$ to $(n,e)$. 

\end{lem}

\begin{proof}
    Consider the sequence of group homomorphisms
    \begin{equation*}
        G \ast_H (H \times \Z) \to [G/K_i \ast_{H_i/K_i} (H_i/K_i \times \Z)] \times G
    \end{equation*}
    defined by sending $g\in G$ to $(gK_i, g)$ and $n\in \Z$ to $(n,e)$. It is easy to check that this sequence of group homomorphisms is trace-preserving in the limit, so it induces a trace-preserving embedding,
    \begin{equation}\label{eqn: group alg map}
    \begin{split}
        \C[G \ast_H (H \times \Z)] \hookrightarrow &\prod_\cU C_r^\ast([G/K_i \ast_{H_i/K_i} (H_i/K_i \times \Z)] \times G)\\
        = &\prod_\cU C_r^\ast(G/K_i \ast_{H_i/K_i} (H_i/K_i \times \Z)) \otimes C_r^\ast(G).
    \end{split}
    \end{equation}

    It suffices to show this embedding is continuous under the reduced norm. By Fell's absorption principle, this is indeed isometric on $\C[G]$ and thus induces an embedding,
    \begin{equation}\label{eqn: embedding}
        C_r^\ast(G) \hookrightarrow \prod_\cU C_r^\ast(G/K_i \ast_{H_i/K_i} (H_i/K_i\times\Z)) \otimes C_r^\ast(G).
    \end{equation}

    Let $E:C_r^\ast(G) \to C_r^\ast(H)$ and $E_i: C_r^\ast(G/K_i) \to C_r^\ast(H_i/K_i)$ be the respective conditional expectations. Let $q_i: \C[G] \to C_r^\ast(G/K_i)$ be the map induced from the quotient map $G \to G/K_i$. Identifying the image of the embedding above with $C_r^\ast(G)$, it is easy to check that the conditional expectation $E$ under this identification sends $(q_i(\lambda_g) \otimes \lambda_g)_\cU$ to $(q_i(E(\lambda_g)) \otimes \lambda_g)_\cU$ for $g \in G$.

    Fix a trace-preserving embedding $C_r^\ast(\Z) \subset C[-1, 1]$. Then,
    \begin{equation}\label{eqn: group embedding}
    \begin{split}
        &C_r^\ast(G/K_i \ast_{H_i/K_i} (H_i/K_i\times\Z)) \otimes C_r^\ast(G)\\
        \subset &[C_r^\ast(G/K_i) \ast_{C_r^\ast(H_i/K_i)} (C_r^\ast(H_i/K_i) \otimes C[-1,1])] \otimes C_r^\ast(G).
    \end{split}
    \end{equation}

    Per Theorem \ref{isomorphism thm}, we furthermore have $C_r^\ast(G/K_i) \ast_{C_r^\ast(H_i/K_i)} (C_r^\ast(H_i/K_i) \otimes C[-1,1]) \subset \cT(E_i)$, so,
    \begin{equation}\label{eqn: embedding 2}
    \begin{split}
        &\prod_\cU C_r^\ast(G/K_i \ast_{H_i/K_i} (H_i/K_i\times\Z)) \otimes C_r^\ast(G)\\
        \subset &\prod_\cU [C_r^\ast(G/K_i) \ast_{C_r^\ast(H_i/K_i)} (C_r^\ast(H_i/K_i) \otimes C[-1,1])] \otimes C_r^\ast(G)\\
        \subset &\prod_\cU \cT(E_i) \otimes C_r^\ast(G),
    \end{split}
    \end{equation}
    where we note that the embedding acts as the identity map on the tensor component $C_r^\ast(G)$, restricts to the natural embedding $C_r^\ast(G/K_i) \hookrightarrow \cT(E_i)$, and sends the semicircular generator of $C[-1, 1]$ to $\frac{1}{2}(T_i + T_i^\ast)$ in $\cT(E_i)$.
    
    Now note that, for $g \in G$,
    \begin{equation*}
        (T_i^\ast q_i(\lambda_g)T_i \otimes \lambda_g)_\cU = (E_i(q_i(\lambda_g)) \otimes \lambda_g)_\cU = (q_i(E(\lambda_g)) \otimes \lambda_g)_\cU
    \end{equation*}
    where the latter equality can be verified easily using the fact that $H_i$ decreases to $H$. Indeed, if $g \in H$, then $q_i(\lambda_g) = \lambda_{gK_i} \in C_r^\ast(H_i/K_i)$, so
    \begin{equation*}
        E_i(q_i(\lambda_g)) \otimes \lambda_g = q_i(\lambda_g) \otimes \lambda_g = q_i(E(\lambda_g)) \otimes \lambda_g.
    \end{equation*}

    If, on the other hand, $g \notin H$, then for large enough $i$, $g \notin H_i$. Thus, $gK_i \notin H_i/K_i$. Hence, again for large enough $i$,
    \begin{equation*}
        E_i(q_i(\lambda_g)) \otimes \lambda_g = 0 = q_i(E(\lambda_g)) \otimes \lambda_g.
    \end{equation*}
    
    This proves the claimed equality. So, by taking linear combinations and norm limits, this implies the copy of $C_r^\ast(G)$ in $\prod_\cU \cT(E_i) \otimes C_r^\ast(G)$, via the embeddings in Equations (\ref{eqn: embedding}) and (\ref{eqn: embedding 2}), together with the isometry $(T_i)_\cU$, satisfies the universal property defining $\cT(E)$. Hence, there is a $\ast$-homomorphism
    \begin{equation*}
        \cT(E) \to \prod_\cU \cT(E_i) \otimes C_r^*(G)
    \end{equation*}
    which extends the embedding of $C_r^\ast(G)$ into $\prod_\cU \cT(E_i) \otimes C_r^\ast(G)$ and sends $T$ to $(T_i)_\cU$. By Theorem \ref{isomorphism thm}, restricting to the algebra generated by $C_r^\ast(G)$ and $\frac{1}{2}(T + T^\ast)$, we then have a $\ast$-homomorphism,
    \begin{equation}\label{eqn: semicircular map}
    \begin{split}
        &C_r^\ast(G) \ast_{C_r^\ast(H)} (C_r^\ast(H) \otimes C[-1,1])\\
        \to &\prod_\cU [C_r^\ast(G/K_i) \ast_{C_r^\ast(H_i/K_i)} (C_r^\ast(H_i/K_i) \otimes C[-1,1])] \otimes C_r^\ast(G).
    \end{split}
    \end{equation}

    This map sends $\lambda_g \in C_r^\ast(G)$ to $\lambda_{gK_i} \otimes \lambda_g$ and restricts to the identity map on $C[-1, 1]$, as it sends $\frac{1}{2}(T + T^\ast)$ to $\left(\frac{1}{2}(T_i + T_i^\ast)\right)_\cU$.

    Again, using the previously fixed trace-preserving embedding $C_r^\ast(\Z) \subset C[-1, 1]$ as in Equation (\ref{eqn: group embedding}), we can restrict to a $\ast$-homomorphism,
    \begin{equation*}
    \begin{split}
        &C_r^\ast(G \ast_H (H \times \Z))\\
        = &C_r^\ast(G) \ast_{C_r^\ast(H)} (C_r^\ast(H) \otimes C_r^\ast(\Z))\\
        \hookrightarrow &\prod_\cU [C_r^\ast(G/K_i) \ast_{C_r^\ast(H_i/K_i)} (C_r^\ast(H_i/K_i) \otimes C_r^\ast(\Z))] \otimes C_r^\ast(G)\\
        = &\prod_\cU C_r^\ast([G/K_i \ast_{H_i/K_i} (H_i/K_i \times \Z)] \times G).
    \end{split}
    \end{equation*}

    Since the map in Equation (\ref{eqn: semicircular map}) acts as the identity map on $C[-1, 1]$, the map above acts as the identity map on $\Z$. Moreover, since the map in Equation (\ref{eqn: semicircular map}) sends $\lambda_g \in C_r^\ast(G)$ to $\lambda_{gK_i} \otimes \lambda_g$, the map above sends $g \in G$ to $(gK_i, g) \in [G/K_i \ast_{H_i/K_i} (H_i/K_i \times \Z)] \times G$. This means the map above is exactly the extension of the embedding in Equation (\ref{eqn: group alg map}), which proves that the embedding there is indeed continuous under the reduced norm, as claimed.
\end{proof}

\end{document}
